\documentclass[10pt,a4paper]{amsart}
\usepackage[margin=19mm]{geometry}
\usepackage{amsmath,amssymb,mathtools}
\usepackage[scr=rsfs,cal=euler]{mathalfa}
\usepackage{xfrac}
\usepackage[unicode,hidelinks,bookmarks=false]{hyperref}
\newcommand{\ie}{i.e.\ }
\newcommand{\cf}{cf.\ }
\newcommand{\resp}{respectively\ }
\newcommand{\Subsection}{\subsection}
\providecommand{\nobreakdash}{\nobreak\hbox{-}}
\setlength{\emergencystretch}{2em}
\multlinegap0pt
\usepackage{bm}
\usepackage{bbm}
\usepackage{dsfont}
\usepackage{xspace}
\usepackage{cancel}
\usepackage{mathtools}
\usepackage{amsthm}
\makeatletter
\@ifundefined{theorem}{\newtheorem{theorem}{Theorem}}{}
\@ifundefined{lemma}{\newtheorem{lemma}[theorem]{Lemma}}{}
\makeatother
\newcommand{\bv}{\mathbf{v}}
\newcommand{\mdso}{\ddot{\mathcal{S}}}
\newcommand{\trans}{^\top}
\title[Inverse eigenvalue problem for discrete Schr\"odinger operat]{Inverse eigenvalue problem for discrete Schr\"odinger operators of a graph}
\author{Anzila Laikhuram, Jephian C.-H. Lin}
\date{Source: arXiv:2506.15430v2 (CC BY 4.0)}
\begin{document}
\maketitle

\noindent\emph{Source and licence.} Inverse eigenvalue problem for discrete Schr\"odinger operators of a graph. Version arXiv:2506.15430v2 (CC BY 4.0), distributed under the Creative Commons Attribution 4.0 International licence, as recorded in the arXiv licence metadata for this exact version and shown on the abstract page https://arxiv.org/abs/2506.15430v2. Full rights notice: licenses/arxiv-2506.15430-CC-BY-4.0.txt.\par\smallskip

\noindent\textit{Editorial scope.} Counted unit: the Lemma whose source label is \texttt{lem:givens} in the article (source0.tex lines 528--531, source label \texttt{lem:givens}), delivered together with the article's own complete original proof of it (source0.tex lines 532--574, one \texttt{proof} environment of 43 lines). No mathematics is written, repaired or re-derived: statement and proof are the article's own text line for line. Required context retained uncounted: none. Every definition, display and label of those lines is retained unchanged. Omitted from this package and not redistributed: 1--527, 575--1217. Nothing inside a retained excerpt is omitted or altered: every retained body is delivered exactly as the article's lines are written, so no omission receipt of any kind accompanies this package. Citations: the retained excerpts carry no citation command at all, so no bibliography entry of the article is transcribed and no reference is redistributed. Reference adaptation: none was needed and none was made; every reference inside the delivered unit resolves inside it, so no unresolved reference or citation is produced. Printed slips kept literal: no printed word, symbol, hypothesis or number of the counted statement or of its proof was altered. Layout and class: the article's own class is \texttt{article} with packages including inputenc, fontenc, amsmath, amssymb, amsthm, hyperref, enumitem, cite, cleveref, array, tikz, soul, cancel; the wrapper is the lane's pinned \texttt{amsart} editorial wrapper at 10pt on A4, which restates the theorem-like environments the retained text needs and adds the article's own macro definitions that the retained text needs (\texttt{\textbackslash bv}, \texttt{\textbackslash mdso}, \texttt{\textbackslash trans}), with extra wrapper packages (bm, bbm, dsfont, xspace, cancel, mathtools). The delivered numbering is the unit's own. Assets: no retained span contains a figure environment, a graphics-inclusion command, a PDF-page inclusion, a table or a TikZ picture; no image file is copied into this package, the package declares no asset dependency and no image file is redistributed with it. Licence: version arXiv:2506.15430v2 of this article is distributed under the Creative Commons Attribution 4.0 International licence, as recorded in the arXiv licence metadata for this exact version and shown on the abstract page https://arxiv.org/abs/2506.15430v2. A full rights notice is delivered at licenses/arxiv-2506.15430-CC-BY-4.0.txt. Characters: every retained excerpt is pure ASCII, so the delivered engine, XeLaTeX with its default Latin Modern text font, reports no missing character.
\begin{lemma}
\label{lem:givens}
Let $G$ be a graph with two vertices $i$ and $j$ such that $N(i)\setminus\{j\} \supseteq N(j)\setminus\{i\}$.  Define $H$ as the graph obtained from $G$ by adding edges from $j$ to each vertex in $N(i)\cup\{i\}\setminus\{j\}$ if the edge was not present.  If $A = \begin{bmatrix} a_{i,j} \end{bmatrix} \in \mdso(G)$ such that either $i \sim j$ in $G$ or $a_{i,i} < a_{j,j}$ when $i \nsim j$ in $G$, then the matrix $A' = R_{i,j;\theta} A R_{i,j;\theta}\trans$ is a matrix in $\mdso(H)$ that is similar to $A$ for any $\theta > 0$ small enough.
\end{lemma}

\begin{proof}
Without loss of generality, we may assume $i = 1$ and $j = 2$.  Let $n$ be the order of $G$.  Thus, we may write 
\[
    A = \begin{bmatrix}
        A_{1,1} & A_{1,2} \\
        A_{1,2}\trans & A_{2,2}
    \end{bmatrix}
    \text{ and }
    R_{i,j;\theta} = \begin{bmatrix}
        R_\theta & O_{2,n-2} \\
        O_{n-2,2} & I_{n-2}
    \end{bmatrix}
\]
such that $A_{1,1}$ is of order $2$.  Consequently, 
\[
    R_{i,j;\theta} A R_{i,j;\theta}\trans = \begin{bmatrix}
        R_\theta A_{1,1} R_\theta\trans & R_\theta A_{1,2} \\
        A_{1,2}\trans R_\theta\trans & A_{2,2}
    \end{bmatrix}.
\]
Since $N(i)\setminus\{j\} \supseteq N(j)\setminus\{i\}$, each nonzero column $\bv$ in $A_{1,2}$ is of the form
\[
    \begin{bmatrix} - \\ - \end{bmatrix}
    \text{ or } 
    \begin{bmatrix} - \\ 0 \end{bmatrix},
\] where each ``$-$'' represents a negative real number.
Thus, for $\theta$ small enough, $R_\theta \bv$ is a entrywise negative vector.  

Suppose $i \sim j$.  Then the fact that $A_{1,1}\in\mdso(K_2)$ implies $R_\theta A_{1,1} R_\theta\trans$ is in $\mdso(K_2)$ when $\theta$ is small enough.  Suppose $i \nsim j$ and $a_{i,i} < a_{j,j}$.  Then by direct computation
\[
    \begin{aligned}
    R_\theta A_{1,1} R_\theta\trans &= 
    \begin{bmatrix} \cos\theta & -\sin\theta \\ \sin\theta & \cos\theta \end{bmatrix}
    \begin{bmatrix} a_{i,i} & 0 \\ 0 & a_{j,j} \end{bmatrix}
    \begin{bmatrix} \cos\theta & \sin\theta \\ -\sin\theta & \cos\theta \end{bmatrix} \\
    &= \begin{bmatrix}
        a_{i,i}\cos^2\theta + a_{j,j}\sin^2\theta & (a_{i,i} - a_{j,j})\sin\theta\cos\theta \\
        (a_{i,i} - a_{j,j})\sin\theta\cos\theta & a_{i,i}\sin^2\theta + a_{j,j}\cos^2\theta. 
    \end{bmatrix}
    \end{aligned}
\]
With $a_{i,i} < a_{j,j}$ and $\sin\theta\cos\theta = \frac{1}{2}\sin 2\theta$, the matrix $R_\theta A_{1,1} R_\theta\trans$ is in $\mdso(K_2)$ when $\theta > 0$ is small enough.
\end{proof}

\end{document}
